\subsection*{The cost of reading pictured words}
An image model can outperform a token model for reasons other than its input representation. We tested this distinction in two comparisons (Fig.~\ref{fig:dialogue-efficiency}). First, a randomly initialized 817,506-parameter subword Transformer trained on exactly the same dialogue batches as the inherited visual writer. It returned 519/4,413 and 708/5,588 normalized first-reference answers, compared with 1,256 and 1,747 for the writer. The development advantage of 18.59 percentage points had a 95\% whole-story bootstrap interval of 16.97 to 20.19 points. The respective dialogue fits consumed 16.59 and 15.04 kJ. The visual model also had a trained optical reader, passage initialization and a direct word-copy action.

We then trained two 547,125-parameter writers from scratch, changing the letter observation while holding word actions, training data and architecture fixed. The image writer used a measured bitmap for each letter. The byte writer received the letter's discrete identity. The image writer obtained 1,103 reserved and 1,574 development exact answers; the byte writer obtained 1,121 and 1,558. The paired image-minus-byte difference on development was $+0.29$ percentage points (95\% interval $-0.77$ to $+1.35$) and on reserved stories $-0.41$ points ($-1.58$ to $+0.77$). Four-reference development F1 was 41.8\% and 42.2\%. Exchanging the pictured question and preceding exchange reduced image-writer exact answers from 1,574 to 237. Exchanging the passage reduced them to 672.

Both writers first exceeded 20\% exact accuracy on both splits at the 2,400-update checkpoint. Accumulated device energy including preparation was 5.20 kJ for images and 4.86 kJ for bytes. Only the byte writer reached the 25\% threshold on both splits by update 8,000. Neither writer exactly answered any of the 83 questions requiring absent wording in all four human references. This study establishes learning from printed ink at small scale; it provides no evidence for an order-of-magnitude reduction in training energy.

The deterministic word-and-stop criterion correctly identified all physically exact replies in an independent evaluation of the two inherited writers (20,002 assessed replies, zero disagreements). A correct first word can still lead to a wrong phrase. Among the direct writer's 4,367 development replies differing from the reference ink, the first error selected the wrong visible word on 3,404, spelled an incorrect word on 788, stopped early on 60 and failed to stop on 115. On the 697 primary composition cases, 628 first errors selected the wrong visible word. The most frequent exact-image failure therefore occurs when selecting content from the available evidence. Exact-image errors can also reflect acceptable alternate renderings or phrasing.
